\documentclass[12pt,a4paper]{article}

% ==================== PACKAGES ====================
\usepackage[utf8]{inputenc}
\usepackage[vietnamese]{babel}
\usepackage{amsmath,amssymb,amsthm}
\usepackage{geometry}
\usepackage{enumitem}
\usepackage[most]{tcolorbox}
\usepackage{hyperref}
\usepackage{fancyhdr}
\usepackage{graphicx}
\usepackage{tikz}
\usepackage{eso-pic}
\usepackage{xcolor}

\geometry{a4paper, margin=2cm, bottom=2.5cm}
\hypersetup{colorlinks=true, linkcolor=blue!70!black, urlcolor=blue}
\setlist[itemize]{topsep=4pt, itemsep=2pt}
\setlist[enumerate]{topsep=4pt, itemsep=2pt}

% ==================== WATERMARK ====================
\AddToShipoutPictureFG{%
  \AtPageCenter{%
    \begin{tikzpicture}[remember picture, overlay]
      \node[rotate=45, scale=1, opacity=0.06]{%
        \fontsize{60}{72}\selectfont\bfseries\sffamily Đặng Tấn Quí%
      };
    \end{tikzpicture}%
  }%
}

% ==================== HEADER / FOOTER ====================
\pagestyle{fancy}
\fancyhf{}
\fancyhead[L]{\small\textit{Lý thuyết Toán 12}}
\fancyhead[R]{\small\textit{Đặng Tấn Quí}}
\fancyfoot[C]{\thepage}
\fancyfoot[L]{\footnotesize\textcopyright\ Đặng Tấn Quí}
\fancyfoot[R]{\footnotesize SĐT: 0377\,122\,210}
\renewcommand{\headrulewidth}{0.4pt}
\renewcommand{\footrulewidth}{0.4pt}

\fancypagestyle{plain}{%
  \fancyhf{}
  \fancyfoot[C]{\thepage}
  \fancyfoot[L]{\footnotesize\textcopyright\ Đặng Tấn Quí}
  \fancyfoot[R]{\footnotesize SĐT: 0377\,122\,210}
  \renewcommand{\headrulewidth}{0pt}
  \renewcommand{\footrulewidth}{0.4pt}
}

% ==================== CUSTOM ENVIRONMENTS ====================
\newtcolorbox{dinhnghia}[1][]{%
  enhanced, breakable,
  colback=blue!3!white, colframe=blue!60!black,
  fonttitle=\bfseries, title={Định nghĩa}, #1%
}

\newtcolorbox{dinhly}[1][]{%
  enhanced, breakable,
  colback=green!3!white, colframe=green!50!black,
  fonttitle=\bfseries, title={Định lý}, #1%
}

\newtcolorbox{tinhchat}[1][]{%
  enhanced, breakable,
  colback=yellow!5!white, colframe=orange!50!black,
  fonttitle=\bfseries, title={Tính chất}, #1%
}

\newtcolorbox{phuongphap}[1][]{%
  enhanced, breakable,
  colback=gray!5!white, colframe=gray!50!black,
  fonttitle=\bfseries, title={Phương pháp}, #1%
}

\newtcolorbox{luuy}[1][]{%
  enhanced, breakable,
  colback=red!3!white, colframe=red!50!black,
  fonttitle=\bfseries, title={Lưu ý}, #1%
}

\newtcolorbox{congthuc}[1][]{%
  enhanced, breakable,
  colback=cyan!3!white, colframe=cyan!50!black,
  fonttitle=\bfseries, title={Công thức}, #1%
}

\newtcolorbox{vidu}[1][]{%
  enhanced, breakable,
  colback=violet!3!white, colframe=violet!50!black,
  fonttitle=\bfseries, title={Ví dụ}, #1%
}

% ==================== DOCUMENT ====================
\begin{document}

% ==================== TRANG BÌA ====================
\thispagestyle{empty}
\begin{tikzpicture}[remember picture, overlay]
  % Background gradient
  \fill[blue!80!black] (current page.south west) rectangle (current page.north east);
  \fill[blue!60!cyan, opacity=0.8]
    (current page.south west) -- (current page.north west) --
    ([xshift=-3cm]current page.north east) -- (current page.south east) -- cycle;

  % Decorative circles
  \foreach \r/\o in {6/0.05, 4.5/0.08, 3/0.12} {
    \fill[white, opacity=\o] ([xshift=2cm, yshift=-2cm]current page.north east) circle (\r cm);
  }

  % Decorative line
  \draw[white, line width=2pt]
    ([yshift=-5cm]current page.north west) ++(2cm,0) -- ++(17cm,0);
  \draw[white, line width=2pt]
    ([yshift=-16cm]current page.north west) ++(2cm,0) -- ++(17cm,0);

  % Title
  \node[anchor=west, white, font=\fontsize{36}{42}\bfseries\selectfont]
    at ([xshift=3cm, yshift=-7.5cm]current page.north west)
    {LÝ THUYẾT TOÁN 12};

  \node[anchor=west, white, font=\fontsize{18}{24}\selectfont]
    at ([xshift=3cm, yshift=-9.5cm]current page.north west)
    {Tài liệu tổng hợp kiến thức};

  % Math symbols decoration
  \node[white, opacity=0.15, font=\fontsize{120}{120}\selectfont, rotate=-15]
    at ([xshift=-3cm, yshift=4cm]current page.center)
    {$\displaystyle\int$};
  \node[white, opacity=0.12, font=\fontsize{80}{80}\selectfont, rotate=10]
    at ([xshift=5cm, yshift=3cm]current page.center)
    {$\sum$};
  \node[white, opacity=0.10, font=\fontsize{90}{90}\selectfont, rotate=-5]
    at ([xshift=7cm, yshift=-1cm]current page.center)
    {$\pi$};
  \node[white, opacity=0.10, font=\fontsize{70}{70}\selectfont, rotate=20]
    at ([xshift=-5cm, yshift=-3cm]current page.center)
    {$\infty$};

  % Author block
  \node[anchor=west, white, font=\Large\bfseries]
    at ([xshift=3cm, yshift=-13cm]current page.north west)
    {Biên soạn:};
  \node[anchor=west, yellow!80!white, font=\fontsize{28}{34}\bfseries\selectfont]
    at ([xshift=3cm, yshift=-14.5cm]current page.north west)
    {ĐẶNG TẤN QUÍ};

  % Contact
  \node[anchor=west, white!90, font=\large]
    at ([xshift=3cm, yshift=-18cm]current page.north west)
    {SĐT: 0377\,122\,210};

  % Copyright notice at bottom
  \node[anchor=south, white!80, font=\small]
    at ([yshift=1.5cm]current page.south)
    {Tài liệu có bản quyền --- Cấm sao chép, phân phối khi chưa được sự đồng ý của tác giả};

  % Year
  \node[anchor=south east, white, font=\Large\bfseries]
    at ([xshift=-2cm, yshift=3cm]current page.south east)
    {\the\year};
\end{tikzpicture}

\newpage

% ==================== TRANG THÔNG TIN BẢN QUYỀN ====================
\thispagestyle{empty}
\vspace*{5cm}
\begin{center}
  {\Large\bfseries LÝ THUYẾT TOÁN 12}\\[8pt]
  {\large Tài liệu tổng hợp kiến thức}
\end{center}

\vspace{2cm}

\begin{tcolorbox}[
  enhanced, colback=red!3!white, colframe=red!60!black,
  fonttitle=\bfseries, title={Thông tin bản quyền},
  boxrule=1.5pt
]
  \begin{itemize}[leftmargin=*]
    \item \textbf{Tác giả:} Đặng Tấn Quí
    \item \textbf{Liên hệ:} 0377\,122\,210
    \item \textbf{Năm:} \the\year
  \end{itemize}

  \vspace{8pt}
  \textbf{Bản quyền \textcopyright\ \the\year\ Đặng Tấn Quí. Bảo lưu mọi quyền.}

  \vspace{6pt}
  Tài liệu này là tài sản trí tuệ của tác giả. Nghiêm cấm mọi hình thức sao chép, tái bản, phân phối, chia sẻ dưới bất kỳ hình thức nào (in ấn, kỹ thuật số, trực tuyến) mà không có sự đồng ý bằng văn bản của tác giả.

  \vspace{6pt}
  Mọi vi phạm sẽ bị xử lý theo quy định của pháp luật về sở hữu trí tuệ.
\end{tcolorbox}

\vfill
\begin{center}
  \small\textit{Tài liệu lưu hành nội bộ --- Không phát hành thương mại}
\end{center}

\newpage

% ==================== MỤC LỤC ====================
\tableofcontents
\newpage

%% ================================================================
%%  PHẦN I: ĐẠI SỐ & GIẢI TÍCH
%% ================================================================
\section{Đại số và Giải tích}

\subsection{Ứng dụng đạo hàm để khảo sát và vẽ đồ thị hàm số}

% ------ 1.1.1 Tính đơn điệu ------
\subsubsection{Tính đơn điệu của hàm số}

\begin{dinhnghia}
  \begin{itemize}
    \item Hàm số $y = f(x)$ gọi là \textbf{đồng biến (tăng)} trên $K$ nếu
      \[
        \forall\, x_1, x_2 \in K,\quad x_1 < x_2 \;\Rightarrow\; f(x_1) < f(x_2).
      \]
    \item Hàm số $y = f(x)$ gọi là \textbf{nghịch biến (giảm)} trên $K$ nếu
      \[
        \forall\, x_1, x_2 \in K,\quad x_1 < x_2 \;\Rightarrow\; f(x_1) > f(x_2).
      \]
    \item Hàm số đồng biến hoặc nghịch biến gọi chung là \textbf{đơn điệu}.
  \end{itemize}
\end{dinhnghia}

\begin{luuy}
  \begin{itemize}
    \item Nếu hàm số đồng biến trên $K$ thì từ trái sang phải đồ thị đi lên.
    \begin{center}
      \includegraphics[width=0.5\linewidth]{D://DangTanQui//Toan//tai_lieu_toan//cursor//cap_3//lop_12//hinh/hs_dong_bien.png}
    \end{center}
    \item Nếu hàm số nghịch biến trên $K$ thì từ trái sang phải đồ thị đi xuống.
    \begin{center}
      \includegraphics[width=0.5\linewidth]{D://DangTanQui//Toan//tai_lieu_toan//cursor//cap_3//lop_12//hinh/hs_nghich_bien.png}
    \end{center}
  \end{itemize}
\end{luuy}

\begin{dinhly}[title={Điều kiện đủ xét tính đơn điệu}]
  Cho hàm số $f$ có đạo hàm trên khoảng $K$:
  \begin{itemize}
    \item $f'(x) > 0,\;\forall x \in K \;\Rightarrow\; f$ đồng biến trên $K$.
    \item $f'(x) < 0,\;\forall x \in K \;\Rightarrow\; f$ nghịch biến trên $K$.
  \end{itemize}
\end{dinhly}

\begin{dinhly}[title={Tổng quát}]
  Cho hàm số $f$ có đạo hàm trên $K$ ($K$ là khoảng, đoạn hoặc nửa khoảng):
  \begin{itemize}
    \item $f'(x) > 0,\;\forall x \in K$ và $f'(x) = 0$ tại một số hữu hạn điểm trên $K$ $\Rightarrow\; f$ đồng biến trên $K$.
    \item $f'(x) < 0,\;\forall x \in K$ và $f'(x) = 0$ tại một số hữu hạn điểm trên $K$ $\Rightarrow\; f$ nghịch biến trên $K$.
  \end{itemize}
\end{dinhly}

\begin{luuy}
  \begin{itemize}
    \item Nếu hàm số $y = f(x)$ liên tục trên đoạn $[a; b]$ và $f'(x) > 0,\;\forall x \in (a; b)$ thì hàm số đồng biến trên đoạn $[a; b]$.
    \item Nếu hàm số $y = f(x)$ liên tục trên đoạn $[a; b]$ và $f'(x) < 0,\;\forall x \in (a; b)$ thì hàm số nghịch biến trên đoạn $[a; b]$.
    \item Tương tự với các khái niệm hàm số đồng biến, nghịch biến trên các nửa khoảng.
  \end{itemize}
\end{luuy}

\begin{phuongphap}[title={Các bước xét tính đơn điệu}]
  \begin{enumerate}
    \item \textbf{Tìm tập xác định.}
      \begin{itemize}
        \item $\dfrac{1}{P(x)}$ có nghĩa khi $P(x) \ne 0$.
        \item $\sqrt{P(x)}$ có nghĩa khi $P(x) \ge 0$.
        \item $\dfrac{1}{\sqrt{P(x)}}$ có nghĩa khi $P(x) > 0$.
      \end{itemize}
    \item \textbf{Tính} $f'(x)$ và tìm các nghiệm $x_i$ của phương trình $f'(x) = 0$ hoặc tại đó $f'(x)$ không xác định.
    \item \textbf{Lập bảng biến thiên} sắp xếp các điểm $x_i$ theo thứ tự tăng dần và xét dấu $f'(x)$.
    \item \textbf{Kết luận} các khoảng đồng biến, nghịch biến.
  \end{enumerate}
\end{phuongphap}

\begin{vidu}
  Xét tính đơn điệu của hàm số $y = x^3 - 3x + 2$.
\end{vidu}

\medskip\noindent\textbf{Các dạng bài tập:}
\begin{enumerate}[label=\textbf{Dạng \arabic*:}, leftmargin=*, align=left]
  \item Xét tính đơn điệu từ biểu thức.
  \item Xét từ bảng biến thiên hoặc đồ thị.
  \item Tìm tham số $m$ để hàm số đơn điệu trên khoảng cho trước.
\end{enumerate}

% ------ 1.1.2 Cực trị ------
\subsubsection{Cực trị của hàm số}

\begin{dinhnghia}
  \begin{itemize}
    \item $f(x)$ đạt \textbf{cực đại} tại $x_0$ nếu $\exists\, h > 0$ sao cho
      \[
        f(x) < f(x_0),\quad \forall\, x \in (x_0 - h,\; x_0 + h),\; x \ne x_0.
      \]
    \item $f(x)$ đạt \textbf{cực tiểu} tại $x_0$ nếu $\exists\, h > 0$ sao cho
      \[
        f(x) > f(x_0),\quad \forall\, x \in (x_0 - h,\; x_0 + h),\; x \ne x_0.
      \]
    \item Giá trị cực đại, cực tiểu gọi chung là \textbf{cực trị}.
    \begin{center}
      \includegraphics[width=0.5\linewidth]{D://DangTanQui//Toan//tai_lieu_toan//cursor//cap_3//lop_12//hinh/hs_cuc_tri.png}
    \end{center}
  \end{itemize}
\end{dinhnghia}

\begin{dinhly}[title={Điều kiện cần và đủ}]
  Giả sử $f'(x_0) = 0$:
  \begin{itemize}
    \item $f'(x)$ đổi dấu từ $+$ sang $-$ khi $x$ qua $x_0$ $\Rightarrow$ $f$ đạt \textbf{cực đại} tại $x_0$.
    \begin{center}
        \includegraphics[width=0.5\linewidth]{D://DangTanQui//Toan//tai_lieu_toan//cursor//cap_3//lop_12//hinh/cuc_dai.png}
    \end{center}
    \item $f'(x)$ đổi dấu từ $-$ sang $+$ khi $x$ qua $x_0$ $\Rightarrow$ $f$ đạt \textbf{cực tiểu} tại $x_0$.
    \begin{center}
        \includegraphics[width=0.5\linewidth]{D://DangTanQui//Toan//tai_lieu_toan//cursor//cap_3//lop_12//hinh/cuc_tieu.png}
    \end{center}
  \end{itemize}
  \textbf{Quy tắc đạo hàm cấp hai:} Nếu $f'(x_0) = 0$ thì:
  \begin{itemize}
    \item $f''(x_0) > 0 \;\Rightarrow\; x_0$ là điểm cực tiểu.
    \item $f''(x_0) < 0 \;\Rightarrow\; x_0$ là điểm cực đại.
    \item Vậy $f'(x_0) = 0$ và $f''(x_0) \ne 0$ thì hàm số đạt cực trị tại $x_0$.
  \end{itemize}
\end{dinhly}

\begin{luuy}
  \begin{itemize}
    \item Giá trị cực đại, cực tiểu \textbf{không nhất thiết} là giá trị lớn nhất, nhỏ nhất của hàm số.
    \item Nếu $f'(x) = 0$ có $n$ nghiệm đơn $x_i$ thì hàm số có $n$ điểm cực trị. Nếu có nghiệm kép $x_j$ thì điểm $x_j$ không phải là điểm cực trị.
    \item $f'(x_0)$ là hệ số góc của tiếp tuyến tại $x_0$.
    \item \textbf{Phương trình tiếp tuyến} tại điểm $M\bigl(x_0,\, f(x_0)\bigr)$:
      \[
        y = f'(x_0)\,(x - x_0) + f(x_0).
      \]
    \item Tiếp tuyến tại điểm cực trị song song với trục $Ox$ (vì $f'(x_0) = 0$).
  \end{itemize}
\end{luuy}

\begin{phuongphap}[title={Các bước tìm cực trị}]
  \begin{enumerate}
    \item Tìm tập xác định.
    \item \textbf{Tính} $f'(x)$ và tìm các nghiệm $x_i$ của phương trình $f'(x) = 0$ hoặc tại đó $f'(x)$ không xác định.
    \item \textbf{Lập bảng biến thiên} sắp xếp các điểm $x_i$ theo thứ tự tăng dần và xét dấu $f'(x)$ hoặc xét dấu $f''(x)$.
    \item Kết luận cực trị.
  \end{enumerate}
\end{phuongphap}

\begin{vidu}
  Tìm cực trị của hàm số $y = x^3 - 3x^2 + 4x - 2$.
\end{vidu}

\medskip\noindent\textbf{Các dạng bài tập:}
\begin{enumerate}[label=\textbf{Dạng \arabic*:}, leftmargin=*, align=left]
  \item Tìm cực trị của hàm số cho bởi biểu thức.
  \item Cực trị hàm bậc ba.
  \item Cực trị hàm trùng phương.
  \item Cực trị của hàm $y = |f(x)|$, $y = f(|x|)$.
  \item Tìm cực trị khi biết biểu thức $f(x)$ và $f'(x)$.
  \item Tìm $m$ để hàm số đạt cực trị tại $x = x_0$.
    \begin{itemize}
      \item Giải $f'(x_0) = 0$ để tìm $m$.
      \item Kiểm tra $f''(x_0) > 0$ (cực tiểu) hoặc $f''(x_0) < 0$ (cực đại).
    \end{itemize}
  \item Hàm số đa thức bậc cao, hàm căn thức.
  \item Tìm $m$ để hàm số có $n$ cực trị ($f'(x)$ có $n$ nghiệm phân biệt đổi dấu).
  \item Đường thẳng đi qua 2 điểm cực trị.
    \begin{itemize}
      \item Chia $y$ cho $y'$: $y = y' \cdot q(x) + h(x)$.
      \item Tại điểm cực trị $x_0$: $y_0 = h(x_0)$ (vì $y'(x_0) = 0$).
      \item Đường thẳng: $y = h(x)$.
    \end{itemize}
  \item Tìm $m$ để hàm bậc 3 có cực trị thỏa điều kiện cho trước.
  \item Tìm $m$ để hàm trùng phương có cực trị thỏa điều kiện cho trước.
  \item Tìm $m$ để hàm bậc hai trên bậc nhất có cực trị thỏa điều kiện cho trước.
  \item Bài toán cực trị hàm chứa dấu giá trị tuyệt đối.
    \begin{itemize}
      \item $y = |f(x)| = \sqrt{f^2(x)}$.
      \item $y' = \dfrac{2f(x) \cdot f'(x)}{2\sqrt{f^2(x)}} = 0 \;\Rightarrow\; f(x) = 0$ hoặc $f'(x) = 0$.
    \end{itemize}
  \item Số điểm cực trị của hàm hợp.
    \begin{itemize}
      \item $y = f\bigl(g(x)\bigr) \;\Rightarrow\; y' = g'(x) \cdot f'\bigl(g(x)\bigr) = 0$.
      \item $\Rightarrow g'(x) = 0$ hoặc $f'\bigl(g(x)\bigr) = 0$.
    \end{itemize}
  \item Tìm $m$ để hàm $f\bigl(g(x)\bigr)$ thỏa mãn điều kiện cho trước.
  \item Tìm cực trị của $f\bigl(g(x)\bigr)$ hoặc $f\bigl(g(x)\bigr) + h(x)$ khi biết đồ thị $f(x)$ hoặc $f'(x)$.
\end{enumerate}

% ------ 1.1.3 GTLN – GTNN ------
\subsubsection{Giá trị lớn nhất -- nhỏ nhất}

\begin{dinhnghia}
  \begin{itemize}
    \item $\displaystyle\max_{[a,b]} f(x) = M$ nếu $f(x) \le M,\;\forall x \in [a,b]$ và $\exists\, x_0 \in [a,b]$ sao cho $f(x_0) = M$.
    \item $\displaystyle\min_{[a,b]} f(x) = m$ nếu $f(x) \ge m,\;\forall x \in [a,b]$ và $\exists\, x_0 \in [a,b]$ sao cho $f(x_0) = m$.
  \end{itemize}
\end{dinhnghia}

\begin{phuongphap}[title={Tìm GTLN--GTNN trên đoạn $[a,b]$}]
  \begin{enumerate}
    \item Tính $f'(x)$, giải $f'(x) = 0$ trên $(a,b)$ được các nghiệm $x_1, x_2, \ldots$
    \item So sánh các giá trị $f(a),\; f(b),\; f(x_1),\; f(x_2),\; \ldots$
    \item Giá trị lớn nhất trong đó là $\max$, nhỏ nhất là $\min$.
  \end{enumerate}
\end{phuongphap}

\begin{tinhchat}
  \begin{itemize}
    \item Nếu $f$ đồng biến trên $[a,b]$:
      $\;\max f = f(b)$, $\;\min f = f(a)$.
    \item Nếu $f$ nghịch biến trên $[a,b]$:
      $\;\max f = f(a)$, $\;\min f = f(b)$.
  \end{itemize}
\end{tinhchat}

\begin{luuy}
  \begin{itemize}
    \item Nếu hàm số liên tục chỉ có 1 cực đại $x_0$ trên $(a, b)$ thì $\max f = f(x_0)$
    \item Nếu hàm số liên tục chỉ có 1 cực tiểu $x_0$ trên $(a, b)$ thì $\min f = f(x_0)$
  \end{itemize}
\end{luuy}

\begin{congthuc}[title={Bất đẳng thức thường dùng}]
  \begin{itemize}
    \item \textbf{AM--GM (Cô-si):} Với $a, b \ge 0$,
      \[
        a + b \ge 2\sqrt{ab}, \qquad \text{dấu ``$=$'' xảy ra} \;\Leftrightarrow\; a = b.
      \]
    \item \textbf{Bunhiacopxki (Cauchy--Schwarz):}
      \[
        |ax + by| \le \sqrt{(a^2 + b^2)(x^2 + y^2)}, \qquad \text{dấu ``$=$'' xảy ra} \;\Leftrightarrow\; \dfrac{a}{x} = \dfrac{b}{y}.
      \]
  \end{itemize}
\end{congthuc}

\medskip\noindent\textbf{Các dạng bài tập:}
\begin{enumerate}[label=\textbf{Dạng \arabic*:}, leftmargin=*, align=left]
  \item Tìm GTLN--GTNN trên đoạn bằng công thức, bảng biến thiên, đồ thị.
  \item Tìm GTLN--GTNN bằng phương pháp đổi biến.
  \item Bài toán chứa tham số $m$.
  \item Đặt ẩn phụ tìm điều kiện $m$ để $f(x,m) = 0$ có nghiệm.
  \item Đặt ẩn phụ tìm $m$ để bất phương trình có nghiệm hoặc nghiệm đúng $\forall x \in K$.
  \item Bài toán thực tế (AM--GM, Cauchy--Schwarz).
  \item GTLN--GTNN qua đồ thị, bảng biến thiên.
  \item GTLN--GTNN trên đoạn.
  \item GTLN--GTNN trên khoảng (tính $\lim_{x \to a^+}$, $\lim_{x \to b^-}$).
  \item Tìm $m$ để GTLN--GTNN thỏa điều kiện.
  \item Tìm $m$ với hàm chứa dấu giá trị tuyệt đối.
  \item GTLN--GTNN hàm ẩn, hàm hợp.
  \item Ứng dụng GTLN--GTNN giải bài toán thực tế.
\end{enumerate}

% ------ 1.1.4 Đường tiệm cận ------
\subsubsection{Đường tiệm cận}

\begin{dinhnghia}
  \begin{itemize}
    \item \textbf{Tiệm cận ngang} $y = L$\; nếu $\displaystyle\lim_{x \to \pm\infty} f(x) = L$.
    \item \textbf{Tiệm cận đứng} $x = a$\; nếu $\displaystyle\lim_{x \to a^{\pm}} f(x) = \pm\infty$.
    \item \textbf{Tiệm cận xiên} $y = kx + l$ với
      \[
        k = \lim_{x \to \pm\infty} \dfrac{f(x)}{x}, \qquad
        l = \lim_{x \to \pm\infty} \bigl[f(x) - kx\bigr].
      \]
  \end{itemize}
\end{dinhnghia}

\begin{congthuc}[title={Hàm phân thức $y = \dfrac{ax + b}{cx + d}$ \;$(c \ne 0,\; ad - bc \ne 0)$}]
  \[
    \text{TCĐ:}\; x = -\dfrac{d}{c}, \qquad
    \text{TCN:}\; y = \dfrac{a}{c}.
  \]
\end{congthuc}

\medskip\noindent\textbf{Các dạng bài tập:}
\begin{enumerate}[label=\textbf{Dạng \arabic*:}, leftmargin=*, align=left]
  \item Tìm tiệm cận của đồ thị hàm số cho bởi công thức.
  \item Tìm tiệm cận biết bảng biến thiên, đồ thị hoặc hàm số liên quan.
  \item Tiệm cận của hàm hợp.
  \item Bài toán tiệm cận chứa tham số $m$.
  \item Xác định tiệm cận thông qua bảng biến thiên, đồ thị.
  \item Xác định tiệm cận thông qua hàm số cho trước.
\end{enumerate}

% ------ 1.2 Khảo sát hàm số ------
\subsection{Khảo sát sự biến thiên và vẽ đồ thị hàm số}

\begin{phuongphap}[title={Sơ đồ khảo sát hàm số}]
  \begin{enumerate}
    \item Tìm tập xác định.
    \item Tính đạo hàm, giải $f'(x) = 0$ (hoặc tìm điểm $f'$ không tồn tại).
    \item Tìm các giới hạn tại $\pm\infty$ và tại các điểm gián đoạn để xác định tiệm cận.
    \item Lập bảng biến thiên (xét dấu $f'(x)$, tìm cực trị).
    \item Tìm các điểm đặc biệt: giao với $Ox$, $Oy$, tâm đối xứng.
    \item Vẽ đồ thị.
  \end{enumerate}
\end{phuongphap}

\begin{tinhchat}[title={Các dạng hàm số thường khảo sát}]
  \begin{itemize}
    \item Hàm bậc ba: $y = ax^3 + bx^2 + cx + d$.
    \item Hàm trùng phương: $y = ax^4 + bx^2 + c$.
    \item Hàm nhất biến: $y = \dfrac{ax + b}{cx + d}$.
    \item Hàm phân thức bậc hai trên bậc nhất: $y = \dfrac{ax^2 + bx + c}{px + q}$.
  \end{itemize}
\end{tinhchat}

\begin{luuy}
  \textbf{Ứng dụng đồ thị:} Dùng đồ thị biện luận số nghiệm phương trình $f(x) = m$.
\end{luuy}

\begin{vidu}
  \begin{itemize}
    \item Khảo sát và vẽ đồ thị hàm số $y = x^3 - 3x^2 + 2$.
    \item Dựa vào đồ thị $y = x^3 - 3x$, tìm $m$ để $x^3 - 3x = m$ có 3 nghiệm phân biệt.
  \end{itemize}
\end{vidu}

\medskip\noindent\textbf{Các dạng bài tập:}
\begin{enumerate}[label=\textbf{Dạng \arabic*:}, leftmargin=*, align=left]
  \item Nhận dạng hàm số thông qua đồ thị (bậc 3, trùng phương, nhất biến).
  \item Xét dấu hệ số hàm số thông qua đồ thị.
  \item Đồ thị hàm chứa dấu giá trị tuyệt đối.
  \item Tương giao đồ thị thông qua đồ thị, bảng biến thiên.
  \item Tương giao đồ thị thông qua hàm số cho trước.
  \item Tương giao đường thẳng với đồ thị hàm bậc 3.
  \item Tương giao đường thẳng với đồ thị hàm nhất biến.
  \item Tương giao đường thẳng với đồ thị hàm trùng phương.
\end{enumerate}

% ------ 1.3 Nguyên hàm ------
\subsection{Nguyên hàm}

\begin{dinhnghia}
  Hàm số $F(x)$ được gọi là \textbf{nguyên hàm} của $f(x)$ trên $K$ nếu $F'(x) = f(x),\;\forall x \in K$.
  \[
    \int f(x)\,dx = F(x) + C \qquad (C \in \mathbb{R}).
  \]
\end{dinhnghia}

\begin{congthuc}[title={Bảng nguyên hàm cơ bản}]
  \begin{align*}
    &\int 0 dx = C
    &\qquad&
    \int k dx = kx + C \\[6pt]
    &\int x^n\,dx = \dfrac{x^{n+1}}{n+1} + C \quad (n \ne -1)
    &\qquad&
    \int \dfrac{1}{x}\,dx = \ln|x| + C \\[6pt]
    &\int \dfrac{1}{\sqrt{x}}\,dx = 2\sqrt{x} + C
    &\qquad&
    \int \dfrac{1}{x^n}\,dx = \dfrac{-1}{(n - 1)x^{n - 1}} + C \\[6pt]
    &\int e^x\,dx = e^x + C
    &&
    \int a^x\,dx = \dfrac{a^x}{\ln a} + C \quad (a > 0,\, a \ne 1) \\[6pt]
    &\int \sin x\,dx = -\cos x + C
    &&
    \int \cos x\,dx = \sin x + C \\[6pt]
    &\int \dfrac{1}{\cos^2 x}\,dx = \tan x + C
    &&
    \int \dfrac{1}{\sin^2 x}\,dx = -\cot x + C \\[6pt]
    &\int \tan x dx = -\ln{|cos x|} + C
    &&
    \int \cot x dx = \ln{|sin x|} + C
  \end{align*}
\end{congthuc}

\begin{congthuc}[title={Một số nguyên hàm bổ sung}]
  \begin{align*}
    &\int \dfrac{f'(x)}{f(x)}\,dx = \ln\bigl|f(x)\bigr| + C
    &\qquad&
    \int \dfrac{dx}{x^2 + a^2} = \dfrac{1}{a}\arctan\dfrac{x}{a} + C \\[6pt]
    &\int \dfrac{dx}{\sqrt{a^2 - x^2}} = \arcsin\dfrac{x}{a} + C
    &&
    \int \dfrac{dx}{x^2 - a^2} = \dfrac{1}{2a}\ln\left|\dfrac{x - a}{x + a}\right| + C
  \end{align*}
\end{congthuc}

\begin{tinhchat}
  \begin{align*}
    &\int f'(x)\,dx = f(x) + C \\[4pt]
    &\int k \cdot f(x)\,dx = k \int f(x)\,dx \\[4pt]
    &\int \bigl[f(x) + g(x)\bigr]\,dx = \int f(x)\,dx + \int g(x)\,dx \\[4pt]
    &\int f(ax + b)\,dx = \dfrac{1}{a} F(ax + b) + C
  \end{align*}
\end{tinhchat}

\begin{congthuc}[title={Nguyên hàm hàm hợp}]
  \[
    \int f\bigl(u(x)\bigr) \cdot u'(x)\,dx = F\bigl(u(x)\bigr) + C.
  \]
\end{congthuc}

\begin{congthuc}[title={Nguyên hàm từng phần}]
  \[
    \int u\,dv = uv - \int v\,du.
  \]
\end{congthuc}

\begin{luuy}
  \begin{itemize}
    \item Nếu $f(x)$ liên tục trên $K$ thì $f(x)$ có nguyên hàm trên $K$.
    \item $f(x)\,dx$ là vi phân của $F(x)$: $dF(x) = F'(x)\,dx = f(x)\,dx$.
  \end{itemize}
\end{luuy}

% ------ 1.4 Tích phân ------
\subsection{Tích phân}

\begin{dinhnghia}
  Nếu $F(x)$ là một nguyên hàm của $f(x)$ trên $[a,b]$ thì:
  \[
    \int_a^b f(x)\,dx = F(x)\bigr|_a^b = F(b) - F(a).
  \]
\end{dinhnghia}

\begin{tinhchat}
  \begin{gather*}
    \int_a^a f(x)\,dx = 0 \\[4pt]
    \int_a^b f(x)\,dx = -\int_b^a f(x)\,dx \\[4pt]
    \int_a^b f(x)\,dx = \int_a^b f(-x)\,dx \\[4pt]
    \int_a^b k \cdot f(x)\,dx = k \int_a^b f(x)\,dx \\[4pt]
    \int_a^b \bigl[f(x) + g(x)\bigr]\,dx = \int_a^b f(x)\,dx + \int_a^b g(x)\,dx \\[4pt]
    \int_a^b f(x)\,dx = \int_a^c f(x)\,dx + \int_c^b f(x)\,dx
  \end{gather*}
  \begin{itemize}
    \item $f(x) \ge 0$ trên $[a,b]$ $\Rightarrow$ $\displaystyle\int_a^b f(x)\,dx \ge 0$.
    \item $f(x) \ge g(x)$ trên $[a,b]$ $\Rightarrow$ $\displaystyle\int_a^b f(x)\,dx \ge \int_a^b g(x)\,dx$.
    \item $m \le f(x) \le M$ trên $[a,b]$ $\Rightarrow$ $m(b-a) \le \displaystyle\int_a^b f(x)\,dx \le M(b-a)$.
    \item $f$ chẵn: $\displaystyle\int_{-a}^{a} f(x)\,dx = 2\int_0^a f(x)\,dx$.
    \item $f$ lẻ: $\displaystyle\int_{-a}^{a} f(x)\,dx = 0$.
  \end{itemize}
\end{tinhchat}

\begin{phuongphap}[title={Các phương pháp tính tích phân}]
  \begin{itemize}
    \item \textbf{Đổi biến:} Đặt $u = g(x)$, đổi cận.
    \item \textbf{Từng phần:}
      \[
        \int_a^b u\,dv = uv\bigr|_a^b - \int_a^b v\,du.
      \]
  \end{itemize}
\end{phuongphap}

\begin{vidu}
  Tính $\displaystyle\int_0^2 (x^2 + 1)\,dx$.
\end{vidu}

\medskip\noindent\textbf{Các dạng bài tập:}
\begin{enumerate}[label=\textbf{Dạng \arabic*:}, leftmargin=*, align=left]
  \item Phân thức hữu tỉ.
  \item Tích của các hàm lượng giác.
\end{enumerate}

% ------ 1.5 Ứng dụng tích phân ------
\subsection{Ứng dụng tích phân}

\begin{congthuc}[title={Diện tích hình phẳng}]
  \textbf{Giữa đồ thị $y = f(x)$ và trục $Ox$:}
  \[
    S = \int_a^b \bigl|f(x)\bigr|\,dx.
  \]
  Nếu $\exists\, c \in [a,b]$: $f(x) \ge 0$ trên $[a,c]$ và $f(x) \le 0$ trên $[c,b]$, thì:
  \[
    S = \int_a^c f(x)\,dx - \int_c^b f(x)\,dx.
  \]
  \textbf{Giữa hai đồ thị $y = f(x)$ và $y = g(x)$:}
  \[
    S = \int_a^b \bigl|f(x) - g(x)\bigr|\,dx.
  \]
\end{congthuc}

\begin{congthuc}[title={Thể tích vật thể tròn xoay}]
  Quay đồ thị $y = f(x)$ quanh trục $Ox$:
  \[
    V = \pi \int_a^b f^2(x)\,dx.
  \]
  \textbf{Giữa hai đồ thị $y = f(x)$ và $y = g(x)$ với $f(x).g(x) \ge 0$:}
  \[
    V = \pi \int_a^b |f^2(x) - g^2(x)|\,dx.
  \]
\end{congthuc}

\medskip\noindent\textbf{Các dạng bài tập:}
\begin{enumerate}[label=\textbf{Dạng \arabic*:}, leftmargin=*, align=left]
  \item Ứng dụng tích phân tìm diện tích.
  \item Ứng dụng tích phân tìm thể tích.
\end{enumerate}

% ------ 1.5 Xác suất có điều kiện ------
\subsection{Xác suất có điều kiện}

\begin{dinhnghia}
  Cho hai biến cố $A$ và $B$ với $P(B) > 0$. \textbf{Xác suất có điều kiện} của $A$ khi biết $B$ đã xảy ra:
  \[
    P(A \mid B) = \dfrac{P(AB)}{P(B)}.
  \]
\end{dinhnghia}

\begin{congthuc}[title={Công thức nhân}]
  \[
    P(AB) = P(B) \cdot P(A \mid B).
  \]
\end{congthuc}

\begin{dinhnghia}[title={Biến cố độc lập}]
  Hai biến cố $A$, $B$ được gọi là \textbf{độc lập} nếu:
  \[
    P(AB) = P(A) \cdot P(B).
  \]
  Khi đó $P(A \mid B) = P(A)$ và $P(B \mid A) = P(B)$.
\end{dinhnghia}

\begin{congthuc}[title={Xác suất toàn phần}]
  \[
    P(A) = P(B) \cdot P(A \mid B) + P(\bar{B}) \cdot P(A \mid \bar{B}).
  \]
\end{congthuc}

\begin{congthuc}[title={Công thức Bayes}]
  \[
    P(B \mid A) = \dfrac{P(B) \cdot P(A \mid B)}{P(A)}.
  \]
\end{congthuc}

\medskip\noindent\textbf{Các dạng bài tập:}
\begin{enumerate}[label=\textbf{Dạng \arabic*:}, leftmargin=*, align=left]
  \item Tính xác suất có điều kiện sử dụng công thức.
  \item Tính xác suất có điều kiện sử dụng sơ đồ hình cây.
  \item Công thức xác suất toàn phần qua bảng thống kê $2 \times 2$ và sơ đồ hình cây.
  \item Công thức Bayes qua bảng thống kê $2 \times 2$ và sơ đồ hình cây.
  \item Kết hợp công thức xác suất toàn phần và Bayes.
  \item Bài toán thực tế -- xác suất toàn phần.
  \item Bài toán thực tế -- công thức Bayes.
\end{enumerate}

%% ================================================================
%%  PHẦN III: HÌNH HỌC
%% ================================================================
\newpage
\section{Hình học không gian tọa độ}

% ------ 3.1 Vectơ ------
\subsection{Vectơ trong không gian}

\begin{dinhnghia}
  \begin{itemize}
    \item \textbf{Vectơ:} đoạn thẳng có hướng, ký hiệu $\overrightarrow{AB}$.
    \item \textbf{Giá} của vectơ: đường thẳng chứa vectơ.
    \item \textbf{Độ dài} (mô-đun): $|\overrightarrow{AB}|$ là khoảng cách từ $A$ đến $B$.
    \item \textbf{Vectơ-không:} điểm đầu trùng điểm cuối, ký hiệu $\vec{0}$.
  \end{itemize}
\end{dinhnghia}

\begin{tinhchat}
  Cho $\vec{a} = (x_1,\, y_1,\, z_1)$, $\vec{b} = (x_2,\, y_2,\, z_2)$:
  \begin{itemize}
    \item \textbf{Cùng phương:} $\vec{a} = k\vec{b}$ $\Leftrightarrow$ $\dfrac{x_1}{x_2} = \dfrac{y_1}{y_2} = \dfrac{z_1}{z_2}$.
    \begin{itemize}
      \item $k > 0$: cùng hướng. \qquad $k < 0$: ngược hướng.
    \end{itemize}
    \item \textbf{Bằng nhau:} cùng hướng và cùng độ dài.
    \item \textbf{Vectơ đối:} $\vec{a} = -\vec{b}$ (cùng độ dài, ngược hướng).
  \end{itemize}
\end{tinhchat}

% ------ 3.2 Phép toán vectơ ------
\subsection{Phép toán vectơ}

\begin{congthuc}[title={Các quy tắc cộng, trừ}]
  \begin{itemize}
    \item \textbf{Quy tắc ba điểm:} $\overrightarrow{AB} + \overrightarrow{BC} = \overrightarrow{AC}$.
    \item \textbf{Quy tắc hình bình hành:} $\overrightarrow{AB} + \overrightarrow{AD} = \overrightarrow{AC}$ \;($ABCD$ là hình bình hành).
    \item \textbf{Quy tắc hình hộp:} $\overrightarrow{AB} + \overrightarrow{AD} + \overrightarrow{AA'} = \overrightarrow{AC'}$ \;($ABCD.A'B'C'D'$ là hình hộp).
    \item \textbf{Phép trừ:} $\overrightarrow{AB} - \overrightarrow{AC} = \overrightarrow{CB}$.
  \end{itemize}
\end{congthuc}

\begin{tinhchat}
  \begin{itemize}
    \item Giao hoán: $\vec{a} + \vec{b} = \vec{b} + \vec{a}$.
    \item Kết hợp: $(\vec{a} + \vec{b}) + \vec{c} = \vec{a} + (\vec{b} + \vec{c})$.
    \item Phần tử trung hòa: $\vec{a} + \vec{0} = \vec{a}$.
    \item Phân phối: $k(\vec{a} + \vec{b}) = k\vec{a} + k\vec{b}$.
  \end{itemize}
\end{tinhchat}

% ------ 3.3 Tọa độ ------
\subsection{Tọa độ trong không gian}

\begin{congthuc}
  Trong hệ tọa độ $Oxyz$ với $\vec{i},\, \vec{j},\, \vec{k}$ là các vectơ đơn vị trên ba trục:
  \begin{itemize}
    \item $\vec{a} = (a_1;\, a_2;\, a_3) = a_1\vec{i} + a_2\vec{j} + a_3\vec{k}$.
    \item $\overrightarrow{AB} = (x_B - x_A;\; y_B - y_A;\; z_B - z_A)$.
    \item \textbf{Độ dài:} $|\overrightarrow{AB}| = \sqrt{(x_B - x_A)^2 + (y_B - y_A)^2 + (z_B - z_A)^2}$.
    \item \textbf{Trung điểm} $M$ của $AB$:
      $\;M = \Bigl(\dfrac{x_A + x_B}{2};\; \dfrac{y_A + y_B}{2};\; \dfrac{z_A + z_B}{2}\Bigr)$.
    \item \textbf{Trọng tâm} tam giác $ABC$:
      $\;G = \Bigl(\dfrac{x_A + x_B + x_C}{3};\; \dfrac{y_A + y_B + y_C}{3};\; \dfrac{z_A + z_B + z_C}{3}\Bigr)$.
    \item \textbf{Nhân vô hướng:} $k\vec{a} = (ka_1,\, ka_2,\, ka_3)$.
  \end{itemize}
\end{congthuc}

\medskip\noindent\textbf{Các dạng bài tập:}
\begin{enumerate}[label=\textbf{Dạng \arabic*:}, leftmargin=*, align=left]
  \item Xác định tọa độ vectơ, tọa độ điểm, đỉnh hình bình hành và hình hộp.
  \item Hình chiếu vuông góc của điểm lên trục và mặt phẳng tọa độ.
  \item Tính độ dài đoạn thẳng.
  \item Xác định tọa độ điểm thỏa điều kiện cho trước.
\end{enumerate}

% ------ 3.4 Tích vô hướng ------
\subsection{Tích vô hướng}

\begin{congthuc}
  \[
    \vec{a} \cdot \vec{b}
    = |\vec{a}| \cdot |\vec{b}| \cdot \cos\bigl(\vec{a},\, \vec{b}\bigr)
    = a_1 b_1 + a_2 b_2 + a_3 b_3.
  \]
  \begin{center}
      \includegraphics[width=0.5\linewidth]{D://DangTanQui//Toan//tai_lieu_toan//cursor//cap_3//lop_12//hinh/tich_vo_huong.png}
  \end{center}
\end{congthuc}

\begin{tinhchat}
  \begin{itemize}
    \item $\vec{a} \cdot \vec{b} = 0 \;\Leftrightarrow\; \vec{a} \perp \vec{b}$ \;(với $\vec{a}, \vec{b} \ne \vec{0}$).
    \item Giao hoán: $\vec{a} \cdot \vec{b} = \vec{b} \cdot \vec{a}$.
    \item Phân phối: $\vec{a} \cdot (\vec{b} + \vec{c}) = \vec{a} \cdot \vec{b} + \vec{a} \cdot \vec{c}$.
    \item $\vec{a} \cdot \vec{a} = |\vec{a}|^2$.
  \end{itemize}
\end{tinhchat}

% ------ 3.5 Tích có hướng ------
\subsection{Tích có hướng}

\begin{congthuc}
  \[
    [\vec{a},\, \vec{b}]
    = \bigl(a_2 b_3 - a_3 b_2;\;\; a_3 b_1 - a_1 b_3;\;\; a_1 b_2 - a_2 b_1\bigr).
  \]
\end{congthuc}

\begin{phuongphap}[title={Cách nhớ nhanh tích có hướng}]
  Viết tọa độ hai vectơ thành 2 hàng, thêm cột đầu vào cuối:
  \[
    \begin{array}{cccc}
      a_1 & a_2 & a_3 & a_1 \\
      b_1 & b_2 & b_3 & b_1
    \end{array}
  \]
  Muốn tìm tọa độ thứ $i$ của $[\vec{a},\, \vec{b}]$, ta \textbf{bỏ cột $i$} rồi lấy \textbf{hai cột liên tiếp còn lại} để lập tích $a_j b_k - a_k b_j$.

  \medskip
  \textbf{Ví dụ:} Cho $\vec{n}_\alpha = (1;\, 1;\, 1)$ và $\vec{n}_\beta = (-1;\, 1;\, 2)$:
  \[
    [\vec{n}_\alpha,\, \vec{n}_\beta]
    = \left(
      \begin{vmatrix} 1 & 1 \\ 1 & 2 \end{vmatrix};\;\;
      \begin{vmatrix} 1 & 1 \\ 2 & -1 \end{vmatrix};\;\;
      \begin{vmatrix} 1 & 1 \\ -1 & 1 \end{vmatrix}
    \right)
    = (1;\; 3;\; 2).
  \]
\end{phuongphap}

\begin{tinhchat}
  \begin{itemize}
    \item $[\vec{a},\, \vec{b}] \perp \vec{a}$ và $[\vec{a},\, \vec{b}] \perp \vec{b}$.
    \begin{center}
        \includegraphics[width=0.5\linewidth]{D://DangTanQui//Toan//tai_lieu_toan//cursor//cap_3//lop_12//hinh/tich_co_huong.png}
    \end{center}
    \item Diện tích hình bình hành: $S_{\text{hbh}} = \bigl|[\vec{a},\, \vec{b}]\bigr| = |\vec{a}|.|\vec{b}|.\sin{(\vec{a}, \vec{b})}$.
    \item Diện tích tam giác: $S_{\triangle} = \dfrac{1}{2}\bigl|[\vec{a},\, \vec{b}]\bigr|$.
    \begin{center}
        \includegraphics[width=0.5\linewidth]{D://DangTanQui//Toan//tai_lieu_toan//cursor//cap_3//lop_12//hinh/tich_co_huong_dien_tich.png}
    \end{center}
  \end{itemize}
\end{tinhchat}

\begin{congthuc}[title={Tích hỗn tạp}]
  \[
    [\vec{a},\, \vec{b}] \cdot \vec{c}
  \]
  \begin{itemize}
    \item Thể tích hình hộp: $V = \bigl|[\overrightarrow{AB},\, \overrightarrow{AD}] \cdot \overrightarrow{AA'}\bigr|$.
    \item Thể tích tứ diện: $V = \dfrac{1}{6}\bigl|[\overrightarrow{AB},\, \overrightarrow{AC}] \cdot \overrightarrow{AD}\bigr|$.
    \item Ba vectơ đồng phẳng $\Leftrightarrow$ $[\vec{a},\, \vec{b}] \cdot \vec{c} = 0$.
    \item Ba đểm A, B, C thẳng hàng $\Leftrightarrow \bigl|\overrightarrow{AB}, \overrightarrow{AC}\bigl| = \vec{0}$
  \end{itemize}
\end{congthuc}

\medskip\noindent\textbf{Các dạng bài tập:}
\begin{enumerate}[label=\textbf{Dạng \arabic*:}, leftmargin=*, align=left]
  \item Xác định vectơ và chứng minh đẳng thức vectơ.
  \item Chứng minh ba điểm thẳng hàng, phân tích vectơ.
  \item Góc giữa hai vectơ, tích vô hướng.
  \item Đẳng thức vectơ.
  \item Phân tích vectơ theo các vectơ cho trước.
\end{enumerate}

% ------ 3.6 Phương trình mặt phẳng ------
\subsection{Phương trình mặt phẳng}

\begin{dinhnghia}
  \begin{itemize}
    \item \textbf{Vectơ pháp tuyến (VTPT):} $\vec{n} = (A;\, B;\, C)$ vuông góc với mặt phẳng.
    \begin{center}
        \includegraphics[width=0.5\linewidth]{D://DangTanQui//Toan//tai_lieu_toan//cursor//cap_3//lop_12//hinh/vector_phap_tuyen_mat_phang.png}
    \end{center}
    \item \textbf{Phương trình tổng quát:} $Ax + By + Cz + D = 0$.
    \item \textbf{Phương trình mặt chắn} (cắt ba trục tại $A(a;0;0)$, $B(0;b;0)$, $C(0;0;c)$):
      \[
        \dfrac{x}{a} + \dfrac{y}{b} + \dfrac{z}{c} = 1.
      \]
    \begin{center}
        \includegraphics[width=0.5\linewidth]{D://DangTanQui//Toan//tai_lieu_toan//cursor//cap_3//lop_12//hinh/mat_chan.png}
    \end{center}
    \item \textbf{Đi qua} $M(x_0;\, y_0;\, z_0)$ với VTPT $\vec{n} = (A;\, B;\, C)$:
      \[
        A(x - x_0) + B(y - y_0) + C(z - z_0) = 0.
      \]
  \end{itemize}
\end{dinhnghia}

\begin{phuongphap}[title={Viết phương trình mặt phẳng}]
  \begin{enumerate}
    \item \textbf{Biết 1 điểm và VTPT:} áp dụng trực tiếp.
    \item \textbf{Biết 1 điểm và 2 VTCP} $\vec{a}$, $\vec{b}$: tìm $\vec{n} = [\vec{a},\, \vec{b}]$, quay về dạng 1.
    \begin{center}
        \includegraphics[width=0.5\linewidth]{D://DangTanQui//Toan//tai_lieu_toan//cursor//cap_3//lop_12//hinh/mat_phang_di_qua_1_diem_va_co_2_vector_chi_phuong.png}
    \end{center}
    \item \textbf{Biết 3 điểm không thẳng hàng:} lập 2 VTCP từ 3 điểm, quay về dạng 2.
    \begin{center}
        \includegraphics[width=0.5\linewidth]{D://DangTanQui//Toan//tai_lieu_toan//cursor//cap_3//lop_12//hinh/mat_phan_di_qua_3_diem.png}
    \end{center}
  \end{enumerate}
\end{phuongphap}

\begin{congthuc}[title={Vị trí tương đối của hai mặt phẳng $(P)$ và $(Q)$}]
  Cho $(P)\!: a_1x + b_1y + c_1z + d_1 = 0$, $(Q)\!: a_2x + b_2y + c_2z + d_2 = 0$.
  \begin{itemize}
    \item \textbf{Trùng nhau:} $\dfrac{a_1}{a_2} = \dfrac{b_1}{b_2} = \dfrac{c_1}{c_2} = \dfrac{d_1}{d_2}$.
    \item \textbf{Song song:} $\dfrac{a_1}{a_2} = \dfrac{b_1}{b_2} = \dfrac{c_1}{c_2} \ne \dfrac{d_1}{d_2}$.
    \item \textbf{Cắt nhau:} $[\vec{n}_1,\, \vec{n}_2] \ne \vec{0}$.
    \begin{center}
        \includegraphics[width=0.5\linewidth]{D://DangTanQui//Toan//tai_lieu_toan//cursor//cap_3//lop_12//hinh/2_mat_phang_cat_nhau.png}
    \end{center}
    \item \textbf{Vuông góc:} $\vec{n}_1 \cdot \vec{n}_2 = 0$, tức $a_1 a_2 + b_1 b_2 + c_1 c_2 = 0$.
    \begin{center}
        \includegraphics[width=0.5\linewidth]{D://DangTanQui//Toan//tai_lieu_toan//cursor//cap_3//lop_12//hinh/2_mat_phang_vuong_goc.png}
    \end{center}
  \end{itemize}
\end{congthuc}

\begin{congthuc}[title={Khoảng cách}]
  \begin{itemize}
    \item \textbf{Điểm $M(x_0;\, y_0;\, z_0)$ đến mặt phẳng $Ax + By + Cz + D = 0$:} chọn $N \in (P)$,
      \[
        d\bigl(M,\, (P)\bigr) = \dfrac{\bigl|\vec{n} \cdot \overrightarrow{MN}\bigr|}{|\vec{n}|} = \dfrac{|Ax_0 + By_0 + Cz_0 + D|}{\sqrt{A^2 + B^2 + C^2}}.
      \]
    \begin{center}
        \includegraphics[width=0.5\linewidth]{D://DangTanQui//Toan//tai_lieu_toan//cursor//cap_3//lop_12//hinh/khoang_cach_1_diem_den_1_mat_phang.png}
    \end{center}
    \item \textbf{Hai mặt phẳng song song:} chọn 1 điểm trên mặt phẳng này, tính khoảng cách đến mặt phẳng kia. Hoặc sử dụng công thức: $d(P, Q) = \dfrac{|D_1 - D_2|}{\sqrt{A^2 + B^2 + C^2}}$. Với:
    \begin{itemize}
        \item $(P)$: $Ax + By + Cz + D_1 = 0$
        \item $(Q)$: $Ax + By + Cz + D_2 = 0$
    \end{itemize}
    \begin{center}
        \includegraphics[width=0.5\linewidth]{D://DangTanQui//Toan//tai_lieu_toan//cursor//cap_3//lop_12//hinh/khoang_cach_giua_2_mat_phang_song_song.png}
    \end{center}
  \end{itemize}
\end{congthuc}

\begin{phuongphap}[title={Hình chiếu và đối xứng qua mặt phẳng}]
  \textbf{Hình chiếu vuông góc} của $M$ lên $(P)$:
  \begin{enumerate}
    \item Viết đường thẳng $d$ qua $M$ có VTCP $= \vec{n}_{(P)}$.
    \item Tìm giao điểm $H$ của $d$ với $(P)$. Điểm $H$ là hình chiếu.
    \begin{center}
        \includegraphics[width=0.5\linewidth]{D://DangTanQui//Toan//tai_lieu_toan//cursor//cap_3//lop_12//hinh/hinh_chieu_vuong_goc_1_diem_len_mat_phang.png}
    \end{center}
  \end{enumerate}
  \textbf{Điểm đối xứng} $M'$ của $M$ qua $(P)$:
  \begin{itemize}
    \item Tìm hình chiếu $H$ (như trên), rồi $M' = 2H - M$.
  \end{itemize}
  \begin{center}
      \includegraphics[width=0.5\linewidth]{D://DangTanQui//Toan//tai_lieu_toan//cursor//cap_3//lop_12//hinh/diem_doi_xung_qua_mat_phang.png}
  \end{center}
\end{phuongphap}

\begin{vidu}
  Viết phương trình mặt phẳng qua $A(1;\, 0;\, 2)$ có VTPT $\vec{n} = (2;\, -1;\, 3)$.
\end{vidu}

\medskip\noindent\textbf{Các dạng bài tập:}
\begin{enumerate}[label=\textbf{Dạng \arabic*:}, leftmargin=*, align=left]
  \item Viết PT mặt phẳng $(P)$ biết 1 điểm và VTPT.
    \begin{itemize}
        \item Mặt phẳng trung trực đoạn $AB$: $M$ là trung điểm $AB$ và $\vec{n_p} = \overrightarrow{AB}$
        \item Qua 1 điểm, song song mặt phẳng $(Q)$ cho trước: $\vec{n_p} = \vec{n_q}$
        \item Qua 1 điểm, vuông góc đường thẳng $(d)$ cho trước: $\vec{n_p} = \vec{u_d}$
        \item Qua 1 điểm, vuông góc đoạn thẳng $AB$ cho trước: $\vec{n_p} = \overrightarrow{AB}$
    \end{itemize}
  \item Viết PT mặt phẳng $(P)$ biết 1 điểm và 2 VTCP.
    \begin{itemize}
        \item Qua $A$, $B$ và song song $(d)$: chọn $A$ hoặc $B$ và $\vec{n_p} = \bigl[\overrightarrow{AB}, \vec{u_d}\bigr]$
        \item Qua $A$, $B$ và song song $CD$: chọn $A$ hoặc $B$ và $\vec{n_p} = \bigl[\overrightarrow{AB}, \overrightarrow{CD}\bigr]$
        \item Chứa $(d)$ và song song $(d')$: chọn $M \in (d)$ và $\vec{n_p} = \bigl[\vec{u_d}, \vec{u_{d'}}\bigr]$
        \item Chứa $(d)$ và song song $AB$: chọn $M \in (d)$ và $\vec{n_p} = \bigl[\overrightarrow{AB}, \vec{u_d}\bigr]$
        \item Qua $A$, $B$ và vuông góc $(Q)$: chọn $A$ hoặc $B$ và $\vec{n_p} = \bigl[\overrightarrow{AB}, \vec{n_q}\bigr]$
        \item Chứa $(d)$ và vuông góc $(Q)$: chọn $M \in (d)$ và $\vec{n_p} = \bigl[\vec{u_d}, \vec{n_q}\bigr]$
        \item Qua 1 điểm và vuông góc với 2 mặt phẳng $(Q)$, $(R)$: $\vec{n_p} = \bigl[\vec{n_q}, \vec{n_r}\bigr]$
        \item Qua 1 điểm và song song với 2 đường thẳng $(d)$, $(d')$: $\vec{n_p} = \bigl[\vec{u_d}, \vec{u_{d'}}\bigr]$
        \item Qua 1 điểm, vuông góc với mặt phẳng $(Q)$ và song song đường thẳng $(d)$: $\vec{n_p} = \bigl[\vec{n_q}, \vec{u_d}\bigr]$
        \item Chứa $(d)$ và đi qua 1 điểm $M \notin (d)$: Lấy $N \in (d)$, $\vec{n_p} = \bigl[\overrightarrow{MN}, \vec{u_d}\bigr]$
        \item Chứa 2 đường thẳng cắt nhau $(d)$ và $(d')$: Lấy $M \in (d)$ hoặc $\in (d')$, $\vec{n_p} = \bigl[\vec{u_d}, \vec{u_{d'}}\bigr]$
        \item Chứa 2 đường thẳng song song $(d)$ và $(d')$: Lấy $M \in (d)$ và $N \in (d')$, $\vec{n_p} = \bigl[\overrightarrow{MN}, \vec{u_d}\bigr]$ hoặc $= \bigl[\overrightarrow{MN}, \vec{u_{d'}}\bigr]$
    \end{itemize}
  \item Viết PT mặt phẳng qua 3 điểm không thẳng hàng.
  \item Viết PT mặt phẳng song song với 1 mặt phẳng, cách khoảng $k$ cho trước.
  \item Viết PT mặt phẳng cắt 3 trục tọa độ tại $A$, $B$, $C$ với $H(a, b, c)$ là trực tâm tam giác $ABC$.
  \[
  (P): ax + by + cz - (a^2 + b^2 + c^2) = 0
  \]
  \item Khoảng cách từ điểm đến mặt phẳng.
  \item Vị trí tương đối hai mặt phẳng.
  \item Hình chiếu vuông góc của điểm lên mặt phẳng.
  \item Điểm đối xứng qua mặt phẳng.
\end{enumerate}

% ------ 3.7 Phương trình đường thẳng ------
\subsection{Phương trình đường thẳng}

\begin{dinhnghia}
  Đường thẳng $d$ qua $M_0(x_0;\, y_0;\, z_0)$ có VTCP $\vec{u} = (a;\, b;\, c)$:
  \begin{itemize}
    \item \textbf{Dạng tham số:} $\;\begin{cases} x = x_0 + at \\ y = y_0 + bt \\ z = z_0 + ct \end{cases}$\quad $(t \in \mathbb{R})$.
    \item \textbf{Dạng chính tắc:} $\;\dfrac{x - x_0}{a} = \dfrac{y - y_0}{b} = \dfrac{z - z_0}{c}$.
    \begin{center}
        \includegraphics[width=0.5\linewidth]{D://DangTanQui//Toan//tai_lieu_toan//cursor//cap_3//lop_12//hinh/duong_thang_va_vector_chi_phong_phap_tuyen.png}
    \end{center}
  \end{itemize}
\end{dinhnghia}

\begin{congthuc}[title={Đường thẳng qua hai điểm $A(x_1;\, y_1;\, z_1)$ và $B(x_2;\, y_2;\, z_2)$}]
  VTCP $\vec{u} = \overrightarrow{AB}$, phương trình chính tắc:
  \[
    \dfrac{x - x_1}{x_2 - x_1} = \dfrac{y - y_1}{y_2 - y_1} = \dfrac{z - z_1}{z_2 - z_1}.
  \]
  \begin{center}
      \includegraphics[width=0.5\linewidth]{D://DangTanQui//Toan//tai_lieu_toan//cursor//cap_3//lop_12//hinh/duong_thang_di_qua_2_diem.png}
  \end{center}
\end{congthuc}

\begin{congthuc}[title={Đường thẳng là giao tuyến của hai mặt phẳng}]
  \[
    d\!:\;
    \begin{cases}
      A_1 x + B_1 y + C_1 z + D_1 = 0 \\
      A_2 x + B_2 y + C_2 z + D_2 = 0
    \end{cases}
    \quad\text{VTCP:}\; \vec{u} = [\vec{n}_1,\, \vec{n}_2].
  \]
  Chọn 1 điểm bằng cách cho $x = x_0$ hoặc $y = y_0$ hoặc $z = z_0$ nào đó (Ví dụ: $0$), rồi giải phương trình tìm 2 tọa độ còn lại
  \begin{center}
      \includegraphics[width=0.5\linewidth]{D://DangTanQui//Toan//tai_lieu_toan//cursor//cap_3//lop_12//hinh/giao_tuyen_2_mat_phang.png}
  \end{center}
\end{congthuc}

\begin{congthuc}[title={Vị trí tương đối giữa hai đường thẳng $d(\vec{u})$ và $d'(\vec{v})$}]
  Chọn $M \in d$, $N \in d'$:
  \begin{itemize}
    \item \textbf{Trùng nhau:}
      $[\vec{u},\, \vec{v}] = \vec{0}$ và $[\vec{u},\, \overrightarrow{MN}] = \vec{0}$.
    \item \textbf{Song song:}
      $[\vec{u},\, \vec{v}] = \vec{0}$ và $[\vec{u},\, \overrightarrow{MN}] \ne \vec{0}$.
    \item \textbf{Cắt nhau:}
      $[\vec{u},\, \vec{v}] \ne \vec{0}$ và $[\vec{u},\, \vec{v}] \cdot \overrightarrow{MN} = 0$.
    \item \textbf{Chéo nhau:}
      $[\vec{u},\, \vec{v}] \ne \vec{0}$ và $[\vec{u},\, \vec{v}] \cdot \overrightarrow{MN} \ne 0$.
    \item \textbf{Vuông góc:}
      $\vec{u} \cdot \vec{v} = 0$.
  \end{itemize}
\end{congthuc}

\begin{congthuc}[title={Vị trí tương đối giữa đường thẳng $d(\vec{u})$ và mặt phẳng $(P)(\vec{n})$}]
  Chọn $M \in d$:
  \begin{itemize}
    \item \textbf{$d$ nằm trên $(P)$:}
      $\vec{u} \cdot \vec{n} = 0$ và $M \in (P)$.
    \item \textbf{$d$ song song $(P)$:}
      $\vec{u} \cdot \vec{n} = 0$ và $M \notin (P)$.
    \item \textbf{$d$ cắt $(P)$:}
      $\vec{u} \cdot \vec{n} \ne 0$.
    \item \textbf{$d$ vuông góc $(P)$:}
      $[\vec{u},\, \vec{n}] = \vec{0}$ \;(tức $\vec{u} \parallel \vec{n}$).
  \end{itemize}
\end{congthuc}

\begin{phuongphap}[title={Tìm giao điểm}]
  \begin{itemize}
    \item \textbf{Giao điểm hai đường thẳng:} Giải hệ (cho tọa độ bằng nhau), kiểm tra nghiệm.
    \item \textbf{Giao điểm đường thẳng và mặt phẳng:} Thế PT tham số của $d$ vào PT mặt phẳng.
  \end{itemize}
\end{phuongphap}

\begin{congthuc}[title={Khoảng cách}]
  \begin{itemize}
    \item \textbf{Điểm $M$ đến đường thẳng $d$:} chọn $N \in d$,
      \[
        d(M,\, d) = \dfrac{\bigl|[\vec{u},\, \overrightarrow{MN}]\bigr|}{|\vec{u}|}.
      \]
    \item \textbf{Hai đường thẳng chéo nhau $d_1(\vec{u})$ và $d_2(\vec{v})$:} chọn $M \in d_1$, $N \in d_2$,
      \[
        d(d_1,\, d_2) = \dfrac{\bigl|[\vec{u},\, \vec{v}] \cdot \overrightarrow{MN}\bigr|}{\bigl|[\vec{u},\, \vec{v}]\bigr|}.
      \]
  \end{itemize}
\end{congthuc}

\begin{phuongphap}[title={Hình chiếu và đối xứng qua đường thẳng}]
  \textbf{Hình chiếu vuông góc} $H$ của $M$ lên $d$:
  \begin{itemize}
    \item Viết $H$ thuộc $d$ (theo tham số $t$), giải $\overrightarrow{MH} \cdot \vec{u} = 0$ tìm $t$.
    \begin{center}
        \includegraphics[width=0.5\linewidth]{D://DangTanQui//Toan//tai_lieu_toan//cursor//cap_3//lop_12//hinh/hinh_chieu_diem_len_duong_thang.png}
    \end{center}
  \end{itemize}
  \textbf{Điểm đối xứng:} $M' = 2H - M$.
  \begin{center}
      \includegraphics[width=0.5\linewidth]{D://DangTanQui//Toan//tai_lieu_toan//cursor//cap_3//lop_12//hinh/diem_doi_xung_qua_duong_thang.png}
  \end{center}
\end{phuongphap}

\begin{vidu}
  \begin{itemize}
    \item Tìm giao điểm đường thẳng
      $\begin{cases} x = 1 + t \\ y = 2 - t \\ z = 3 + 2t \end{cases}$
      với mặt phẳng $x + y + z - 6 = 0$.
    \item Tính khoảng cách từ $A(1;\, 2;\, 3)$ đến mặt phẳng $2x - y + 2z - 6 = 0$.
  \end{itemize}
\end{vidu}

\medskip\noindent\textbf{Các dạng bài tập:}
\begin{enumerate}[label=\textbf{Dạng \arabic*:}, leftmargin=*, align=left]
  \item Xác định VTCP của đường thẳng.
  \item Viết phương trình đường thẳng $(d)$ đi qua 1 điểm và có VTCP.
  \begin{itemize}
      \item Qua 2 điểm $A$, $B$: Chọn $A$ hoặc $B$ và $\vec{u_d} = \overrightarrow{AB}$
      \item Qua $A$ và song song $(d')$: $\vec{u_d} = \vec{u_{d'}}$
      \item Qua $A$ và vuông góc mặt phẳng $(P)$: $\vec{u_d} = \vec{n_p}$
  \end{itemize}
  \item Viết phương trình đường thẳng $(d)$ đi qua 1 điểm, có 2 VTPT.
  \begin{itemize}
      \item Qua $A$ và vuông góc với 2 đường thẳng $(d_1)$ và $(d_2)$: $\vec{u_d} = \bigl[\vec{u_{d_1}}, \vec{u_{d_2}}\bigr]$
      \item Qua $A$ và song song với 2 mặt phẳng $(P)$ và $(Q)$ (hoặc chứa trong 1 mặt phẳng): $\vec{u_d} = \bigl[\vec{n_p}, \vec{n_q}\bigr]$
      \item Giao tuyến của 2 mặt phẳng $(P)$ và $(Q)$: Tìm $I \in (P) \cap (Q)$ và $\vec{u_d} = \bigl[\vec{n_p}, \vec{n_q}\bigr]$
      \item Qua $A$ và vuông góc với đường thẳng $(d')$ và song song với mặt phẳng $(P)$: $\vec{u_d} = \bigl[\vec{u_{d'}}, \vec{n_p}\bigr]$
      \item Qua $A$ và vuông góc với đường thẳng $(d_1)$ và cắt đường thẳng $(d_2)$: $\vec{u_d} = \bigl[\vec{u_{d_1}}, \vec{n_p}\bigr]$ với $(P)$ là mặt phẳng đi qua $A$ và $(d_2)$
  \end{itemize}
  \item Qua $A$ và vuông góc và cắt đường thẳng $(d')$: $\vec{u_d} = \overrightarrow{AB}$ với $B$ là hình chiếu của $A$ lên $(d')$. Hoặc chọn $M \in (d')$ và tìm pháp tuyến mặt phẳng $(P)$ qua $A$ và chứa $(d')$ là $\vec{n_p} = \bigl[\vec{u_{d'}}, \overrightarrow{AM}\bigr]$. Khi đó: $\vec{u_d} = \bigl[\vec{n_p}, \vec{u_{d'}}\bigr] = \bigl[\bigl[\vec{u_{d'}}, \overrightarrow{AM}\bigr], \vec{u_{d'}}\bigr]$.
  \item Là hình chiếu của đường thẳng $(d')$ lên mặt phẳng $(P)$: Lấy $A, B \in (d')$ tìm hình chiếu $A', B'$ lên mặt phẳng $(P)$ khi đó $\vec{u_d} = \overrightarrow{A'B'}$.
  \item Qua $A$ và cắt 2 đường thẳng $(d_1)$ và $(d_2)$: Lấy $M \in (d_1)$ và $N \in (d_2)$, tìm pháp tuyến mặt phẳng $(P)$ qua $A$ và chứa $(d_1)$ là $\vec{n_p} = \bigl[\vec{u_{d_1}}, \overrightarrow{AM}\bigr]$ và mặt phẳng $(Q)$ qua $A$ và chứa $(d_2)$ là $\vec{n_q} = \bigl[\vec{u_{d_2}}, \overrightarrow{AN}\bigr]$. Khi đó: $\vec{u_d} = \bigl[\vec{n_p}, \vec{n_q}\bigr] = \bigl[\bigl[\vec{u_{d_1}}, \overrightarrow{AM}\bigr], \bigl[\vec{u_{d_2}}, \overrightarrow{AN}\bigr]\bigr]$.
  \item Qua $A$, cắt đường thẳng $(d')$ và song song mặt phẳng $(P)$: Lấy $M \in (d')$, tìm pháp tuyến mặt phẳng $(Q)$ qua $A$ và chứa $(d')$ là $\vec{n_q} = \bigl[\vec{u_{d'}}, \overrightarrow{AM}\bigr]$. Khi đó: $\vec{u_d} = \bigl[\vec{n_p}, \vec{n_q}\bigr] = \bigl[\vec{n_p}, \bigl[\vec{u_{d'}}, \overrightarrow{AM}\bigr]\bigr]$
  \item Là đường vuông góc chung của 2 đường thẳng chéo nhau $(d_1)$ và $(d_2)$: $\vec{u_d} = \bigl[\vec{u_{d_1}}, \vec{u_{d_2}}\bigr]$. Lấy $M \in (d_1)$, tìm mặt phẳng $(P)$ qua $M$ có pháp tuyến $\vec{n_p} = \bigl[\vec{u_d}, \vec{u_{d_1}}\bigr]$, tìm giao điểm $A = (d_2) \cap (P)$.
  \item Vị trí tương đối (đường thẳng - đường thẳng, đường thẳng - mặt phẳng).
  \item Hình chiếu vuông góc của điểm lên đường thẳng.
  \item Điểm đối xứng qua đường thẳng.
\end{enumerate}

% ------ 3.8 Góc trong không gian ------
\subsection{Góc trong không gian}

\begin{congthuc}[title={Góc giữa hai đường thẳng $d(\vec{u})$ và $d'(\vec{v})$}]
  \[
    \cos\bigl(d,\, d'\bigr)
    = \dfrac{|\vec{u} \cdot \vec{v}|}{|\vec{u}| \cdot |\vec{v}|}
    = \dfrac{|aa' + bb' + cc'|}{\sqrt{a^2 + b^2 + c^2}\;\sqrt{a'^2 + b'^2 + c'^2}}.
  \]
  \begin{center}
      \includegraphics[width=1\linewidth]{D://DangTanQui//Toan//tai_lieu_toan//cursor//cap_3//lop_12//hinh/goc_giua_2_duong_thang.png}
  \end{center}
\end{congthuc}

\begin{congthuc}[title={Góc giữa đường thẳng $d(\vec{u})$ và mặt phẳng $(P)(\vec{n})$}]
  \[
    \sin\bigl(d,\, (P)\bigr)
    = \dfrac{|\vec{u} \cdot \vec{n}|}{|\vec{u}| \cdot |\vec{n}|}
    = \dfrac{|aA + bB + cC|}{\sqrt{a^2 + b^2 + c^2}\;\sqrt{A^2 + B^2 + C^2}}.
  \]
  \begin{center}
      \includegraphics[width=0.5\linewidth]{D://DangTanQui//Toan//tai_lieu_toan//cursor//cap_3//lop_12//hinh/goc_giua_duong_thang_va_mat_phang.png}
  \end{center}
\end{congthuc}

\begin{congthuc}[title={Góc giữa hai mặt phẳng $(P)(\vec{n}_1)$ và $(Q)(\vec{n}_2)$}]
  \[
    \cos\bigl((P),\, (Q)\bigr)
    = \dfrac{|\vec{n}_1 \cdot \vec{n}_2|}{|\vec{n}_1| \cdot |\vec{n}_2|}
    = \dfrac{|A_1 A_2 + B_1 B_2 + C_1 C_2|}{\sqrt{A_1^2 + B_1^2 + C_1^2}\;\sqrt{A_2^2 + B_2^2 + C_2^2}}.
  \]
  \begin{center}
      \includegraphics[width=0.5\linewidth]{D://DangTanQui//Toan//tai_lieu_toan//cursor//cap_3//lop_12//hinh/gioc_giua_2_mat_phang.png}
  \end{center}
\end{congthuc}

\medskip\noindent\textbf{Các dạng bài tập:}
\begin{enumerate}[label=\textbf{Dạng \arabic*:}, leftmargin=*, align=left]
  \item Góc giữa hai đường thẳng.
  \item Góc giữa đường thẳng và mặt phẳng.
  \item Góc giữa hai mặt phẳng.
\end{enumerate}

% ------ 3.9 Phương trình mặt cầu ------
\subsection{Phương trình mặt cầu}

\begin{dinhnghia}
  \begin{itemize}
    \item \textbf{Dạng chính tắc:}
      $(x - a)^2 + (y - b)^2 + (z - c)^2 = R^2$,
      \;tâm $I(a;\, b;\, c)$, bán kính $R$.
    \item \textbf{Dạng khai triển:}
      $x^2 + y^2 + z^2 - 2ax - 2by - 2cz + d = 0$
      \;(điều kiện $a^2 + b^2 + c^2 - d > 0$),
      \;tâm $I(a;\, b;\, c)$, bán kính $R = \sqrt{a^2 + b^2 + c^2 - d}$.
     \begin{center}
         \includegraphics[width=0.5\linewidth]{D://DangTanQui//Toan//tai_lieu_toan//cursor//cap_3//lop_12//hinh/mat_cau.png}
     \end{center}
  \end{itemize}
\end{dinhnghia}

\begin{phuongphap}[title={Viết phương trình mặt cầu}]
  \begin{itemize}
    \item \textbf{Biết tâm và bán kính:} áp dụng trực tiếp.
    \item \textbf{Biết tâm và 1 điểm:} $R = |IM|$.
    \item \textbf{Qua 4 điểm:} giải hệ $IA = IB = IC = ID$ tìm tâm.
    \item \textbf{Tiếp xúc mặt phẳng:} mặt cầu tâm $I(a;\,b;\,c)$ tiếp xúc mặt phẳng $(P)$ thì $R = d\bigl(I,\,(P)\bigr)$.
    \item \textbf{Tiếp xúc đường thẳng:} mặt cầu tâm $I$ tiếp xúc đường thẳng $\Delta$ thì $R = d(I,\,\Delta)$.
  \end{itemize}
\end{phuongphap}

\begin{congthuc}[title={Phương trình mặt phẳng tiếp xúc mặt cầu}]
  Mặt phẳng tiếp xúc với mặt cầu $S(I,\,R)$ tại điểm $M_0(x_0;\,y_0;\,z_0)$ trên mặt cầu có VTPT $\vec{n} = \overrightarrow{IM_0}$:
  \[
    (x_0 - a)(x - x_0) + (y_0 - b)(y - y_0) + (z_0 - c)(z - z_0) = 0.
  \]
\end{congthuc}

\begin{congthuc}[title={Vị trí tương đối}]
  \textbf{Điểm $M$ và mặt cầu $S(I, R)$:}
  \begin{itemize}
    \item $|IM| < R$: $M$ nằm trong. \quad
      $|IM| = R$: $M$ nằm trên. \quad
      $|IM| > R$: $M$ nằm ngoài.
  \end{itemize}

  \textbf{Mặt phẳng $(P)$ và mặt cầu $S(I, R)$:}
  Gọi $d = d\bigl(I,\, (P)\bigr) = \dfrac{|Aa + Bb + Cc + D|}{\sqrt{A^2 + B^2 + C^2}}$:
  \begin{itemize}
    \item $d > R$: không giao nhau.
    \item $d = R$: tiếp xúc tại 1 điểm.
    \item $d < R$: cắt nhau theo đường tròn bán kính $r = \sqrt{R^2 - d^2}$.
    \begin{center}
        \includegraphics[width=1\linewidth]{D://DangTanQui//Toan//tai_lieu_toan//cursor//cap_3//lop_12//hinh/vivi_tri_mat_phang_va_mat_cau.png}
    \end{center}
  \end{itemize}

  \textbf{Đường thẳng $\Delta$ và mặt cầu $S(I, R)$:}
  Chọn $M \in \Delta$, gọi $d = d(I,\, \Delta) = \dfrac{\bigl|[\vec{u},\, \overrightarrow{MI}]\bigr|}{|\vec{u}|}$:
  \begin{itemize}
    \item $d > R$: không cắt nhau.
    \item $d = R$: tiếp xúc tại 1 điểm.
    \item $d < R$: cắt tại 2 điểm.
    \begin{center}
        \includegraphics[width=1\linewidth]{D://DangTanQui//Toan//tai_lieu_toan//cursor//cap_3//lop_12//hinh/vi_tri_duong_thang_va_mat_cau.png}
    \end{center}
  \end{itemize}
\end{congthuc}

\begin{vidu}
  Cho mặt cầu tâm $I(1;\, 2;\, 3)$, $R = 5$. Viết phương trình mặt cầu. Tìm giao với mặt phẳng $x = 1$.
\end{vidu}

\medskip\noindent\textbf{Các dạng bài tập:}
\begin{enumerate}[label=\textbf{Dạng \arabic*:}, leftmargin=*, align=left]
  \item Xác định tâm, bán kính, nhận biết phương trình mặt cầu.
  \item Mặt cầu có tâm và đi qua một điểm.
  \item Mặt cầu có đường kính (tìm trung điểm $\to$ tâm).
  \item Mặt cầu qua 4 điểm không đồng phẳng ($IA = IB = IC = ID$).
  \item Mặt cầu có tâm thuộc đường thẳng / mặt phẳng.
  \item Mặt cầu tiếp xúc đường thẳng / mặt phẳng.
  \item Bài toán thực tế.
\end{enumerate}

\end{document}
